%% kapa: exam=1 round=29 year=2018 period=1 form=A subject=real_estate subject_ko=부동산학원론 seq=18 number=98
%% source: Q-Net 감정평가사 기출문제 (cst00309) artlSeq=5202808 file=extracted/1008/5202808/2018년 제29회 감정평가사 1차시험 기출문제/2018년 제29회 감정평가사 1차시험 기출문제/1교시 A형.pdf
%% conversion: kordoc 4.24.51 -> tools/gen1.py (auto); review=proofread-sonnet; flags=none
\begin{question}{98}
한국주택금융공사법령에 의한 주택담보노후연금제도에 관한 설명으로 옳지 않은 것은?
\begin{choices}
  \item 주택소유자와 그 배우자 모두 60세 이상이어야 이용할 수 있다.
  \item 연금지급방식으로 주택소유자가 선택하는 일정한 기간 동안 노후생활자금을 매월 지급받는 방식이 가능하다.
  \item 주택담보노후연금보증을 받은 사람은 담보주택의 소유권등기에 한국주택금융공사의 동의 없이는 제한물권을 설정하거나 압류 등의 목적물이 될 수 없는 재산임을 부기등기 하여야 한다.
  \item 주택담보노후연금을 받을 권리는 양도하거나 압류할 수 없다.
  \item 한국주택금융공사는 주택담보노후연금보증을 받으려는 사람에게 소유주택에 대한 저당권 설정에 관한 사항을 설명하여야 한다.
\end{choices}
\end{question}
\answer{1}
